%% chapters/chapter01/chapter01.tex
%% CHAPTER 1: INTRODUCTION

\chapterfigpath{chapter01}

\chapter{Introduction}
\label{chap:introduction}


\section{Background}
\label{sec:background}

Provide background information about your research domain. Explain the broader context and why this area is important. You might discuss the evolution of the field, current trends, or real-world applications.

Continue with more context. Introduce key concepts that readers need to understand your research. Use citations to support your statements \cite{sample_ref1}.

Figure~\ref{fig:research_framework} illustrates the overall research framework adopted in this study.

\begin{figure}[ht]
    \centering
    \includegraphics[width=0.7\textwidth]{research_framework}
    \caption{Research framework showing the relationship between research components}
    \label{fig:research_framework}
\end{figure}


\section{Problem Statement}
\label{sec:problem_statement}

Despite advances in this field, several challenges remain:

\begin{enumerate}
    \item \textbf{Challenge 1:} Describe the first problem or limitation
    
    \item \textbf{Challenge 2:} Describe the second problem or limitation
    
    \item \textbf{Challenge 3:} Describe the third problem or limitation
\end{enumerate}

Explain how these challenges affect users, researchers, or industry and why solving them is important.


\section{Research Objectives}
\label{sec:objectives}

The primary objective of this research is to state your main goal. The specific objectives are:

\begin{enumerate}
    \item To analyze existing approaches to the problem domain
    
    \item To design and implement a system that addresses identified gaps
    
    \item To evaluate the performance of the proposed solution
    
    \item To provide recommendations for practitioners and researchers
\end{enumerate}


\section{Research Questions}
\label{sec:research_questions}

This research addresses the following questions:

\begin{enumerate}
    \item RQ1: Your first research question - should be answerable through your research
    
    \item RQ2: Your second research question
    
    \item RQ3: Your third research question
\end{enumerate}


\section{Scope and Limitations}
\label{sec:scope}

\subsection{Scope}
\label{subsec:scope}

This research focuses on:
\begin{itemize}
    \item What technologies/platforms you're using
    \item What specific problem aspects you're addressing
    \item What user groups or scenarios you're targeting
    \item What time period or data you're considering
\end{itemize}

\subsection{Limitations}
\label{subsec:limitations}

The study has the following limitations:
\begin{itemize}
    \item Limitation 1 - e.g., ``Limited to English language data''
    \item Limitation 2 - e.g., ``Testing performed on simulated environment only''
    \item Limitation 3 - e.g., ``Sample size of N participants''
    \item Limitation 4 - e.g., ``Does not address real-time scenarios''
\end{itemize}


\section{Significance of the Study}
\label{sec:significance}

This research contributes to the field by:

\begin{itemize}
    \item Contribution 1 - e.g., ``Providing a novel approach to...''
    \item Contribution 2 - e.g., ``Offering practical guidelines for...''
    \item Contribution 3 - e.g., ``Creating a reusable framework for...''
    \item Contribution 4 - e.g., ``Filling a gap in existing literature on...''
\end{itemize}


\section{Thesis Organization}
\label{sec:organization}

This thesis is organized as follows:

\textbf{Chapter 2: Literature Review} examines existing research on the topic, including key theoretical foundations and related work.

\textbf{Chapter 3: Methodology} describes the research design, tools used, and implementation details.

\textbf{Chapter 4: Results and Analysis} presents experimental findings and discusses their implications.

\textbf{Chapter 5: Conclusion} summarizes contributions, discusses limitations, and suggests future research directions.
